\section*{Chapter \ref{ch:rc:curve-construction} --- Curve Construction}

\begin{solution}{exo:rc:curve-construction:1}
Flat forwards equal the discrete forward on the interval: $(3.30\times3 -
3.25\times2)/(3-2) = 3.40\%$, at 2.5 years as everywhere between the \omterm{def:rc:curve-construction:pillar}{pillars}.
\end{solution}

\begin{solution}{exo:rc:curve-construction:2}
The futures rate is $100-96.72 = 3.28\%$; less the convexity adjustment,
$3.28-0.006 = 3.274\%$, simple ACT/360 over the quarter.
\end{solution}

\begin{solution}{exo:rc:curve-construction:3}
Its period is already priced by the deposits and the first futures: two
instruments would price one period. The curve implies a one-year swap rate of
$3.4639\%$; put back at $3.4839\%$, the system is still square (its \omterm{def:rc:curve-construction:pillar}{pillar}
falls between the third and fourth futures) and calibrates exactly, but the
only freedom left is the two weeks between 15 and 29 September 2027: under
flat forwards the forward there jumps from $3.242\%$ to $3.753\%$, and falls
to $3.150\%$ for the rest of the fourth future's quarter. Exactness is bought
with a 51-basis-point spike that every instrument over that fortnight inherits.
\end{solution}

\begin{solution}{exo:rc:curve-construction:4}
$f = z + Tz'$. Left slope $(3.2883-3.4165)/(2.0027-1.2110) = -0.1619\%$ a year:
$f^- = 3.2883 + 2.0027\times(-0.1619) = 2.964\%$. Right slope
$(3.3095-3.2883)/1 = 0.0212\%$: $f^+ = 3.2883+2.0027\times0.0212 = 3.331\%$. A
37-basis-point jump made by the rule.
\end{solution}

\begin{solution}{exo:rc:curve-construction:5}
$f(0,10) = \beta_0 + \beta_1e^{-10/\tau_1} + \beta_2(10/\tau_1)e^{-10/\tau_1} +
\beta_3(10/\tau_2)e^{-10/\tau_2} = 3.9407\%$. As $T\to0$ each loading
$(1-e^{-x})/x\to1$ and the hump terms vanish: $z(0)=\beta_0+\beta_1 =
2.3264\%$.
\end{solution}

\begin{solution}{exo:rc:curve-construction:6}
The four-, five- and twelve-year hedges (and the small three- and fifteen-year
ones): under flat forwards the eight-year date depends only on the seven- and
ten-year \omterm{def:rc:curve-construction:pillar}{pillars}. The spline is a global fit: moving a \omterm{def:rc:curve-construction:pillar}{pillar} bends the whole
zero curve, so the eight-year discount factor responds to \omterm{def:rc:curve-construction:pillar}{pillars} far away,
with alternating signs; those hedges hedged the interpolation, not the market.
\end{solution}

\begin{solution}{exo:rc:curve-construction:7}
With one-basis-point bumps the monotone convex buckets are $-5\,663$ (five
years), $+50\,916$ (seven), $+33\,188$ (ten) and $-7\,232$ (twelve), summing to
USD~71\,216; with a hundredth of a basis point, scaled, they sum to 69\,610,
and a parallel one-basis-point bump gives 69\,579. The scheme's case depends on
the inputs, so a full-basis-point bump of one \omterm{def:rc:curve-construction:pillar}{pillar} switches cases on its
intervals and the finite difference picks up more than the local slope; the
small bump measures the derivative, and its buckets add up.
\end{solution}

\begin{solution}{exo:rc:curve-construction:8}
Repricing the inputs is necessary, not sufficient: every interpolation of the
chapter does it, and they disagree on every instrument between the \omterm{def:rc:curve-construction:pillar}{pillars}
(forwards by tens of basis points, \omterm{def:rc:curve-construction:pillar}{off-pillar} swaps by thousands of dollars).
The curve must also be judged on its forwards (no sawtooth, no overshoot), its
locality and the hedges it implies, and on out-of-sample instruments such as
\omterm{def:rc:curve-construction:pillar}{off-pillar} swaps traded in the market.
\end{solution}

\begin{solution}{pb:rc:curve-construction:1}
\textbf{1.} Each was traded at the par rate of the Friday curve.
\textbf{2.} No: every input still reprices exactly under the new rule, since
calibration forces it.
\textbf{3.} The new curve puts higher forwards between 15 and 20 years (and so
lower ones elsewhere in that span of \omterm{def:rc:curve-construction:pillar}{pillars}): the 17-year par rate rises by
about one basis point, $3.9975\%$ to $4.0074\%$.
\textbf{4.} Every \omterm{def:rc:curve-construction:pillar}{off-pillar} swap, here all seven; a swap maturing on a \omterm{def:rc:curve-construction:pillar}{pillar}
would be unaffected in value.
\textbf{5.} $3.5005\%$ then $3.5007\%$: six years sits in the five-to-seven-year
interval, where both rules give nearly the same average forward.
\textbf{6.} $-\text{USD}~90\,062$.
\textbf{7.} A loss. The 17-year receiver loses USD~863\,596 and the 22-year payer
gains 678\,267; the 13-year payer gains 234\,745; the 11-year receiver, 9-year
payer and 8-year receiver roughly offset.
\textbf{8.} It is a valuation change, not a trading result: product control and
the valuation committee decide how a methodology change is reported, often as a
separate line, and the change is approved before it goes live.
\textbf{9.} ``The new interpolation, which the Treasury also uses, values our
\omterm{def:rc:curve-construction:pillar}{off-pillar} swaps USD~90 thousand lower; no market moved.''
\textbf{10.} No: pillar-maturity swaps are priced by their own quotes whatever the
interpolation.
\textbf{11.} Pay fixed on USD~0.3016 billion of the fifteen-year before, 0.5076
billion after: $346\,584/1\,149\,075$ and $583\,281/1\,149\,075$ billion.
\textbf{12.} Receive fixed on USD~0.2490 billion before, 0.3350 billion after.
\textbf{13.} The sum is the book's parallel DV01, fixed by the level of the curve,
on which both rules agree; the rules disagree only on how the risk is spread
between \omterm{def:rc:curve-construction:pillar}{pillars}.
\textbf{14.} Under monotone convex a full-basis-point bump can switch the case on
an interval and add a spurious piece; a small bump, scaled, measures the
derivative (exercise~7).
\textbf{15.} The forward shape between ten and fifteen years now depends on the
twelve-year \omterm{def:rc:curve-construction:pillar}{pillar} through the \omterm{def:rc:curve-construction:pillar}{pillar} forwards the method estimates; the 11-
and 13-year swaps read that shape.
\textbf{16.} For: continuous, non-overshooting forwards, and the method of the
official Treasury curve, so marks align with a public reference. Against: its
buckets depend on the bump size and it is harder to explain.
\textbf{17.} As a model change: documented, independently validated, its P\&L
impact measured in parallel runs before the switch, approved, and reported
separately.
\textbf{18.} A reserve for the spread of values across acceptable
interpolations (here at least USD~90 thousand on this book), released as the
book's \omterm{def:rc:curve-construction:pillar}{off-pillar} positions run off.
\textbf{19.} Named result: \emph{the interpolation switch} costs
$-\text{USD}~90\,062$ with no market move, and the fifteen-year hedge grows by
68\%, from USD~0.30 to 0.51 billion.
\textbf{20.} Everything between the quotes: the forwards, the value of every
\omterm{def:rc:curve-construction:pillar}{off-pillar} instrument, and where its hedges sit.
\end{solution}

\begin{solution}{iq:rc:curve-construction:1}
Curve bootstrapping solves one \omterm{def:rc:curve-construction:pillar}{pillar} at a time, which works when each
instrument depends only on \omterm{def:rc:curve-construction:pillar}{pillars} up to its own maturity, as with linear or
\omterm{def:rc:curve-construction:flatfwd}{flat-forward interpolation}. A global fit solves all \omterm{def:rc:curve-construction:pillar}{pillars} together (Newton
on the full system); it is needed for non-local interpolations (splines,
tension, monotone convex with its \omterm{def:rc:curve-construction:pillar}{pillar} forwards), for instruments that
depend on later \omterm{def:rc:curve-construction:pillar}{pillars}, or for least-squares fits with more instruments than
\omterm{def:rc:curve-construction:pillar}{pillars}.
\iqlookfor{the lower-triangular structure as the condition, and when it fails.}
\end{solution}

\begin{solution}{iq:rc:curve-construction:2}
The forward is $z+Tz'$: the slope of the zero curve jumps at each \omterm{def:rc:curve-construction:pillar}{pillar}, and
the jump is multiplied by maturity, so the forward curve is a sawtooth with
steps of tens of basis points at long \omterm{def:rc:curve-construction:pillar}{pillars}, which then price forward-starting
instruments and caps badly.
\iqlookfor{the formula $f = z + Tz'$.}
\end{solution}

\begin{solution}{iq:rc:curve-construction:3}
Exact fit of inputs; positive, continuous and smooth forwards; no overshoot;
locality of a bump; stability (small input change, small curve change);
linearity in the inputs (additive buckets); speed; simple explanation.
Smoothness conflicts with locality (a spline is smooth and global); no
overshoot conflicts with linearity (monotone convex is nonlinear).
\iqlookfor{a list and at least one real trade-off.}
\end{solution}

\begin{solution}{iq:rc:curve-construction:4}
A non-local interpolation such as a spline on zero rates: the risk on distant
\omterm{def:rc:curve-construction:pillar}{pillars} comes from the curve rule, not from the swap's cash flows. Check the
curve's settings, then its locality by bumping one input and plotting the
forward change.
\iqlookfor{diagnosing the interpolation, not the trade.}
\end{solution}

\begin{solution}{iq:rc:curve-construction:5}
A turn: an extra \omterm{def:rc:curve-construction:pillar}{pillar} (or an additive spread) over the few days around the
date, fitted to the futures or deposits that straddle it, so its cost is not
spread over a quarter. Meeting dates: one \omterm{def:rc:curve-construction:pillar}{pillar} per meeting with flat
forwards between them at the front, fitted to overnight-index swaps or
futures, then a smooth rule beyond two or three years.
\iqlookfor{step structures handled explicitly, not smoothed.}
\end{solution}

\begin{solution}{iq:rc:curve-construction:6}
Not necessarily. For a curve linear in its inputs the bucket is independent of
the bump; for a nonlinear rule such as monotone convex, or with a loose solver
tolerance, it is not. Check the solver tolerance first (a bump of a tenth of a
basis point must be well above it), then compare with a small bump scaled up;
if the curve is nonlinear, report the small-bump derivative and state it.
\iqlookfor{nonlinearity versus numerical noise, and how to tell them apart.}
\end{solution}
